\documentclass[11pt,a4paper]{article}

\usepackage{fontspec}
\usepackage[main=russian, english]{babel}
\usepackage{amsmath, amsfonts, amssymb}
\usepackage{graphicx}
\usepackage{booktabs}
\usepackage{cite}
\usepackage{geometry}
\usepackage{xcolor}
\usepackage{url}
\usepackage{hyperref}

\setmainfont{Times New Roman}
\setromanfont{Times New Roman}
\setsansfont{Arial}
\setmonofont{Courier New}

\geometry{margin=2.5cm}

\title{\textbf{Физически обоснованная слепая деконволюция: Адаптивное извлечение секвенограмм через дифференцируемый рендеринг и калибровочную изотропную оптимизацию}}
\author{Ваше Имя\thanks{Автор для корреспонденции: email@example.com} \\ 
        \small \textit{Кафедра вычислительной обработки сигналов}}
\date{\today}

\begin{document}

\maketitle

\begin{abstract}
Акустический мониторинг в сложных средах сталкивается с фундаментальной проблемой перекрывающихся сигналов и многолучевой реверберации. В данной работе представлен новый метод слепого разделения сигналов (Blind Source Separation), извлекающий «секвенограммы» — точные, разреженные векторы временной активации — непосредственно из искаженных спектрограмм. В отличие от традиционной сверточной неотрицательной матричной факторизации (ConvNMF), мы отказываемся от пиксельных шаблонов в пользу дифференцируемого параметрического рендеринга, который архитектурно гарантирует соблюдение законов физики (например, точную привязку обертонов к основному тону). Для решения проблемы «курицы и яйца» при слепой оценке мы предлагаем метод совместной оптимизации сигналов, библиотеки шаблонов на сфере $S^n$ и частотно-зависимого причинного фильтра эха. Устойчивость сложной многокомпонентной системы обеспечивается введением концепции Калибровочной нормировки (Gauge Calibration) и изотропной Евклидовой функции потерь, которая формирует динамическую тригонометрическую маршрутизацию градиентов, избавляя алгоритм от эвристического подбора гиперпараметров.
\end{abstract}

\section{Введение}
Распознавание акустических паттернов (например, эхолокационных криков летучих мышей) в естественных средах осложняется эффектами среды: затуханием и реверберацией. Традиционные методы деконволюции часто терпят неудачу, так как не могут отделить внутреннюю структуру сигнала от внешних искажений среды. Эта задача представляет собой классическую проблему совместной слепой оптимизации: невозможно точно оценить профиль эха без знания идеального шаблона сигнала, а извлечь шаблон невозможно без снятия реверберации.

В данной работе мы представляем физически обоснованную генеративную пайплайн-модель, которая решает эту проблему через совместную (joint) оптимизацию трех компонентов: разреженных векторов активации (\textit{секвенограмм}), физически жестко связанной параметрической библиотеки шаблонов и частотно-зависимого 1D-фильтра эха.

\section{Генеративная модель и Физический рендеринг}

Пусть $S \in \mathbb{R}_{\ge 0}^{F \times T}$ — наблюдаемая спектрограмма. Генеративный процесс $\hat{S} = (H * W) * E$ моделируется следующим образом.

\subsection{Параметрическая библиотека на сфере $S^n$ и дифференцируемый рендеринг}
Традиционные методы ConvNMF оптимизируют шаблоны $W \in \mathbb{R}^{F \times T_w}$ попиксельно, что приводит к размытию и нарушению структуры сигнала (например, независимости движения обертонов от основного тона). Мы предлагаем генерировать шаблоны "на лету" с использованием дифференцируемого Гауссова рендеринга физических кривых.

Каждый канал (класс сигнала) $k$ задается макро-вектором $\vec{n}_k$ на гиперсфере $S^n$. Обучаемая нейросетевая функция $T: S^n \to \mathbb{R}^1 \to \mathbb{R}^2$ преобразует эту точку в набор параметров сплайна: кривую фундаментальной частоты $f_0(t)$ и амплитудные огибающие основного тона $A_1(t)$ и первого обертона $A_2(t)$.
Шаблон $W_k$ формируется аналитически для каждой точки матрицы $(f_{idx}, t_{idx})$:
\begin{equation}
W_k(f_{idx}, t) = A_1(t) \exp\left( - \frac{(f_{idx} - f_0(t))^2}{2\sigma^2} \right) + A_2(t) \exp\left( - \frac{(f_{idx} - 2f_0(t))^2}{2\sigma^2} \right)
\end{equation}
Такая архитектура жестко фиксирует физику распространения гармоник, делая обертон зависимым от $f_0(t)$, и кардинально сокращает размерность пространства поиска.

\subsection{Микро-флуктуации в касательном пространстве}
Для адаптации к естественной вариативности сигналов, форма шаблона не является абсолютно жесткой. В момент времени $t$ сеть может предсказывать вектор микро-флуктуаций $\vec{v}_k(t) \in \mathbb{R}^{n+1}$. Для сохранения геометрии мы проецируем его на касательное пространство $T_p S^n$:
\begin{equation}
\vec{\Delta}_k(t) = \vec{v}_k(t) - (\vec{v}_k(t) \cdot \vec{n}_k)\vec{n}_k
\end{equation}
Используя Экспоненциальное отображение (Retraction map), мы получаем смещенную точку на сфере: $\vec{n}_k(t) = \vec{n}_k \cos(\|\vec{\Delta}_k\|) + \frac{\vec{\Delta}_k}{\|\vec{\Delta}_k\|} \sin(\|\vec{\Delta}_k\|)$. Это позволяет библиотеке плавно экстраполировать многообразие сигналов, не разрушая базовую структуру класса.

\subsection{Инстанс-сегментация секвенограмм и профиль эха}
Сгенерированные шаблоны сворачиваются по времени с матрицей секвенограмм $H \in \mathbb{R}_{\ge 0}^{K \times T}$. Полученная «чистая» спектрограмма $V$ подвергается реверберации через причинную частотно-зависимую 1D-свертку (Depthwise 1D Convolution) с матрицей эха $E \in \mathbb{R}^{F \times L_e}$, что запрещает нефизичный перенос энергии между частотными бинами.

\section{Изотропная калибровка многозадачной оптимизации}

Совместная оптимизация $H, W$ и $E$ требует наложения строгих структурных ограничений на $H$. Стандартное линейное суммирование штрафов ($\sum \alpha_i \mathcal{L}_i$) требует ручного подбора весов и нестабильно в разреженных пространствах. Мы предлагаем агрегировать компоненты ошибки через \textbf{Евклидову метрику}: $\mathcal{L}_{reg} = \sqrt{\sum \mathcal{L}_i^2}$. 
Однако это корректно только в том случае, если все штрафы находятся в изотропном пространстве. Мы достигаем этого через соблюдение принципа \textit{Единичной локальной возвращающей силы}.

\subsection{Структурные функционалы}
Мы определяем три ключевых физических штрафа для секвенограмм:
\begin{enumerate}
    \item \textbf{Бинаризация (Резкость):} $\mathcal{L}_{sharp} = \sum_t h(t)(1 - h(t))$. Заставляет активации группироваться у 0 или 1.
    \item \textbf{Локализация (Инстанс-сепарация):} Чтобы один канал не пытался описать два разнесенных во времени идентичных сигнала, мы штрафуем временную дисперсию энергии в каждом канале относительно его центра масс $\mu_k$:
    $\mathcal{L}_{loc} = \sum_k \sum_t h_k(t) \cdot (t - \mu_k)^2$. Это заставляет перекрывающиеся сигналы «разбегаться» по разным уровням $k$.
    \item \textbf{Разреженность (Group Lasso):} $\mathcal{L}_{sparse} = \sum_k \max_t (h_k(t))$. Подавляет неиспользуемые каналы в абсолютный нуль.
\end{enumerate}
Дополнительно вводится штраф на микро-флуктуации $\mathcal{L}_{dev} = \sum_t h(t) \|\vec{\Delta}(t)\|^2$, который требует стабильности формы только в моменты звучания сигнала ($h(t)>0$), игнорируя фоновый вакуум.

\subsection{Тригонометрия градиентов}
Откалибровав каждый функционал $\mathcal{L}_i$ так, чтобы амплитуда его вариационной производной локально равнялась единице ($|\frac{\partial \mathcal{L}_i}{\partial h}|_{local} = 1$), мы можем рассмотреть полный градиент Евклидовой суммы. На примере двух штрафов:
\begin{equation}
\frac{\partial \mathcal{L}_{reg}}{\partial h_i} = \frac{\mathcal{L}_1}{\sqrt{\mathcal{L}_1^2 + \mathcal{L}_2^2}} \frac{\partial \mathcal{L}_1}{\partial h_i} + \frac{\mathcal{L}_2}{\sqrt{\mathcal{L}_1^2 + \mathcal{L}_2^2}} \frac{\partial \mathcal{L}_2}{\partial h_i}
\end{equation}
Введя макроскопический угол состояния системы $\theta = \arctan (\mathcal{L}_2 / \mathcal{L}_1)$, уравнение градиентного спуска принимает изящную тригонометрическую форму:
\begin{equation}
\frac{\partial \mathcal{L}_{reg}}{\partial h_i} = \underbrace{\cos(\theta)}_{\text{макро-вес}} \cdot \underbrace{\left(\frac{\partial \mathcal{L}_1}{\partial h_i}\right)}_{\text{микро-градиент}} + \underbrace{\sin(\theta)}_{\text{макро-вес}} \cdot \underbrace{\left(\frac{\partial \mathcal{L}_2}{\partial h_i}\right)}_{\text{микро-градиент}}
\end{equation}

Это уравнение раскрывает двухуровневую систему \textbf{динамической маршрутизации градиентов}:
\begin{itemize}
    \item \textbf{Макроскопический авто-баланс (Угол $\theta$):} Если система успешно подавляет одну ошибку (например, $\mathcal{L}_1 \to 0$), угол $\theta \to \pi/2$. Тогда $\cos(\theta) \to 0$, а $\sin(\theta) \to 1$, и оптимизатор автоматически переключает 100\% вычислительного внимания на оставшуюся проблему. В отличие от жестко заданных гиперпараметров, наша система \textit{динамически вращает вектор градиента} в процессе обучения.
    \item \textbf{Микроскопический баланс:} Внутри скобок находятся откалиброванные локальные силы. Благодаря условию \textit{единичной локальной силы}, в точках конфликта эти градиенты соизмеримы (амплитуды $\approx 1$), что гарантирует, что ни один штраф не раздавит другой из-за разницы физических размерностей.
\end{itemize}

\section{Стратегия совместного обучения}
Поскольку процессы извлечения шаблона и подавления эха неотделимы друг от друга, модель обучается совместно из информированной начальной точки. Секвенограммы $H$ инициализируются локальной RMS-огибающей спектрограммы; профиль эха $E$ — затухающими экспонентами с $\tau$, убывающим с ростом частоты; а библиотека $T$ — вертикальными широкополосными импульсами. В ходе оптимизации эхо забирает на себя дисперсионные хвосты сигнала, библиотека стягивается в идеальные гладкие параметрические кривые, а RMS-огибающие коллапсируют в дельта-пики секвенограмм благодаря калиброванной Евклидовой метрике.

\section{Результаты и обсуждение}
% [TODO: Вставить результаты экспериментов на синтетических и реальных данных. Показать график сходимости угла theta и восстановленные секвенограммы]
Предложенная параметрическая модель кардинально решает проблему структурных искажений, свойственных классической слепой деконволюции. Внедрение дифференцируемого сплайн-рендеринга гарантирует физическую корректность обертонов, а штраф на дисперсию времени позволяет осуществлять инстанс-сегментацию перекрывающихся сигналов одного класса.

Самым значимым теоретическим вкладом работы является доказательство применимости Евклидовой оптимизации для многозадачных сетей при условии локальной градиентной калибровки. Переход от скалярных весов потерь к динамической фазовой тригонометрии ($\cos\theta, \sin\theta$) полностью устраняет необходимость в ручном подборе гиперпараметров, делая архитектуру самобалансирующейся.

\section{Заключение}
Платформа адаптивного извлечения секвенограмм представляет собой мощный дифференцируемый физический движок для акустического анализа. Объединяя геометрию римановых многообразий, физику акустики и теорию оптимизации, данный метод предлагает надежный, свободный от гиперпараметров инструмент для масштабного разбора сложных аудиосцен в условиях сильной реверберации.

\section*{Доступность данных и кода}
Исходный код на Python/PyTorch (включая FFT-свертки и дифференцируемый рендерер) доступен в репозитории: \url{https://github.com/yourusername/AdaptiveSequenograms}.

\bibliographystyle{plain}
\bibliography{references}

\end{document}